\documentclass[11pt]{article}
\usepackage{amsmath,amssymb,amsthm}
\usepackage{geometry}
\geometry{margin=1in}
\newtheorem{theorem}{Theorem}
\newtheorem{lemma}{Lemma}
\newcommand{\prox}{\operatorname{prox}}
\newcommand{\R}{\mathbb{R}}

\begin{document}

\section*{Proximal Gradient Method (Forward–Backward Splitting)}

\section*{Mathematical Formulation}
Let  
\[
F(x):=f(x)+g(x),\qquad x\in\R^{d},
\]
where  

\begin{itemize}
    \item $f:\R^{d}\!\to\!\R$ is convex and $L$–smooth (i.e. differentiable with $L$‑Lipschitz gradient);
    \item $g:\R^{d}\!\to\!\R\cup\{+\infty\}$ is proper, closed and convex (possibly non‑smooth).
\end{itemize}
Given a stepsize $\eta>0$, the \emph{proximal gradient update} is  

\[
\boxed{%
x_{k+1}= \prox_{\eta g}\!\bigl(x_{k}-\eta \nabla f(x_{k})\bigr)},
\qquad k=0,1,2,\dots
\tag{1}
\]

The proximal operator of $g$ is defined by  

\[
\prox_{\eta g}(v):=\arg\min_{u\in\R^{d}}\Bigl\{ g(u)+\frac{1}{2\eta}\|u-v\|^{2}\Bigr\}.
\tag{2}
\]

\section*{Pseudocode}
\begin{verbatim}
Input:   initial point x0, stepsize η ∈ (0,2/L), 
         maxIter K, tolerance ε
Output:  approx. solution xK

x ← x0
for k = 0,…,K-1 do
    // gradient step on the smooth part
    y ← x - η * grad_f(x)
    // proximal step on the nonsmooth part
    x_new ← prox_{η g}(y)
    if ||x_new - x|| ≤ ε then break
    x ← x_new
end for
return x
\end{verbatim}

\section*{Dry Run Example}
Take $g\equiv0$ (so $\prox_{\eta g}= \operatorname{id}$) and  

\[
f(w)=(w-3)^{2},\qquad \nabla f(w)=2(w-3),\qquad \eta=0.1,\qquad w_{0}=0 .
\]

\[
\begin{array}{c|c|c|c}
k & \nabla f(w_{k}) & w_{k+1}=w_{k}-\eta\nabla f(w_{k}) & 
F(w_{k+1})=(w_{k+1}-3)^{2}\\ \hline
0 & 2(0-3)=-6 & w_{1}=0-0.1(-6)=0.6 & (0.6-3)^{2}=5.76\\
1 & 2(0.6-3)=-4.8 & w_{2}=0.6-0.1(-4.8)=1.08 & (1.08-3)^{2}=3.6864\\
2 & 2(1.08-3)=-3.84 & w_{3}=1.08-0.1(-3.84)=1.464 & (1.464-3)^{2}=2.3689
\end{array}
\]

The iterates move monotonically toward the minimiser $w^{\star}=3$ and the objective values decrease.

\newpage
\section*{Assumptions}
\begin{enumerate}
\item \textbf{Smoothness of $f$:} $\nabla f$ is $L$‑Lipschitz,
\[
\|\nabla f(x)-\nabla f(y)\| \le L\|x-y\|,\qquad \forall x,y\in\R^{d}.
\tag{A1}
\]
Equivalently, for all $x,y$,
\[
f(y)\le f(x)+\langle\nabla f(x),y-x\rangle+\frac{L}{2}\|y-x\|^{2}.
\tag{A1'}
\]

\item \textbf{Convexity of $f$ and $g$:} for any $x,y$ and $\theta\in[0,1]$,
\[
f(\theta x+(1-\theta)y)\le \theta f(x)+(1-\theta)f(y),\qquad
g(\theta x+(1-\theta)y)\le \theta g(x)+(1-\theta)g(y).
\tag{A2}
\]

\item \textbf{Existence of a minimiser:} the set 
\[
\mathcal{X}^{\star}:=\arg\min_{x\in\R^{d}}F(x)
\]
is non‑empty and we denote an arbitrary element by $x^{\star}$.

\item \textbf{Stepsize:} $\eta$ satisfies
\[
0<\eta<\frac{2}{L}.
\tag{A3}
\]
The choice $\eta=1/L$ yields the tightest constant in the bound below.
\end{enumerate}

\section*{Main Convergence Theorem}
\begin{theorem}[Ergodic $O(1/k)$ rate]\label{thm:rate}
Let $\{x_{k}\}$ be generated by \eqref{1} with a stepsize $\eta\in(0,2/L)$.  
For any $x^{\star}\in\mathcal{X}^{\star}$,
\[
F(x_{k})-F(x^{\star})\le
\frac{\|x_{0}-x^{\star}\|^{2}}{2\eta\,k},
\qquad\forall k\ge 1.
\tag{3}
\]
Consequently, $F(x_{k})\to F^{\star}:=F(x^{\star})$ and the sequence $\{x_{k}\}$ is bounded.
\end{theorem}

\section*{Theoretical Properties}
\begin{itemize}
\item \textbf{Stability:} The update \eqref{1} is a non‑expansive mapping when $\eta\in(0,2/L)$, guaranteeing that the iterates remain in a bounded region containing $\mathcal{X}^{\star}$.
\item \textbf{Hyperparameter Sensitivity:} The convergence constant in \eqref{3} is proportional to $1/\eta$. Choosing $\eta$ too small slows progress, while $\eta$ close to $2/L$ may violate the descent inequality used in the proof.
\item \textbf{Comparison to Gradient Descent:} When $g\equiv0$, \eqref{1} reduces to standard gradient descent and the same $O(1/k)$ bound holds. For non‑zero $g$, the proximal step allows handling of regularisers such as $\ell_{1}$ norm without sacrificing the rate.
\end{itemize}

\section*{Discussion}
The theorem is sharp for the class of convex $L$‑smooth $f$ and convex $g$. Faster rates (e.g. $O(1/k^{2})$) require additional assumptions such as strong convexity or acceleration (FISTA). The proof hinges on two elementary properties:
\begin{enumerate}
\item the \emph{descent lemma} for $L$‑smooth $f$;
\item the \emph{non‑expansiveness} (more precisely, the \emph{firm non‑expansiveness}) of the proximal operator.
\end{enumerate}
Both are invoked explicitly in the step‑by‑step derivation below.

\newpage
\section*{Preliminaries (Definitions and Lemmas)}
\begin{lemma}[Firm non‑expansiveness of $\prox_{\eta g}$]\label{lem:firm}
For any $u,v\in\R^{d}$,
\[
\|\prox_{\eta g}(u)-\prox_{\eta g}(v)\|^{2}
\le \langle \prox_{\eta g}(u)-\prox_{\eta g}(v),\,u-v\rangle .
\tag{4}
\]
\end{lemma}
\begin{proof}
By definition \eqref{2}, $p=\prox_{\eta g}(u)$ satisfies the optimality condition  
\[
0\in \partial g(p)+\frac{1}{\eta}(p-u).
\]
Similarly for $q=\prox_{\eta g}(v)$. Subtracting the two inclusions and using monotonicity of $\partial g$ yields \eqref{4}. \hfill $\square$
\end{proof}

\begin{lemma}[Descent Lemma]\label{lem:descent}
If $\nabla f$ is $L$‑Lipschitz, then for all $x,y$,
\[
f(y)\le f(x)+\langle\nabla f(x),y-x\rangle+\frac{L}{2}\|y-x\|^{2}.
\tag{5}
\]
\end{lemma}
\begin{proof}
Standard consequence of integrating the Lipschitz condition on $\nabla f$. \hfill $\square$
\end{proof}

\section*{Main Proof (Step‑by‑Step Derivation)}
Fix an arbitrary optimal point $x^{\star}\in\mathcal{X}^{\star}$.  
Define the auxiliary point
\[
v_{k}:=x_{k}-\eta\nabla f(x_{k}).
\tag{6}
\]
By the update rule \eqref{1}, $x_{k+1}= \prox_{\eta g}(v_{k})$.

\bigskip
\noindent\textbf{[Step 1]} (Optimality condition of the proximal step).  
From \eqref{2},
\[
x_{k+1}= \arg\min_{u}\Bigl\{g(u)+\frac{1}{2\eta}\|u-v_{k}\|^{2}\Bigr\}
\]
implies
\[
0\in \partial g(x_{k+1})+\frac{1}{\eta}\bigl(x_{k+1}-v_{k}\bigr).
\tag{7}
\]

\noindent\textbf{[Step 2]} (Variational inequality for $g$).  
For any $u\in\R^{d}$, the monotonicity of $\partial g$ together with \eqref{7} yields
\[
g(u)\ge g(x_{k+1})+\frac{1}{\eta}\langle x_{k+1}-v_{k},\,u-x_{k+1}\rangle .
\tag{8}
\]

\noindent\textbf{[Step 3]} (Apply \eqref{8} with $u=x^{\star}$).  
\[
g(x^{\star})\ge g(x_{k+1})+\frac{1}{\eta}\langle x_{k+1}-v_{k},\,x^{\star}-x_{k+1}\rangle .
\tag{9}
\]

\noindent\textbf{[Step 4]} (Expand the inner product).  
Using $v_{k}=x_{k}-\eta\nabla f(x_{k})$ from \eqref{6},
\[
x_{k+1}-v_{k}=x_{k+1}-x_{k}+\eta\nabla f(x_{k}).
\tag{10}
\]
Insert \eqref{10} into \eqref{9}:
\[
g(x^{\star})\ge g(x_{k+1})+\frac{1}{\eta}
\bigl\langle x_{k+1}-x_{k}+\eta\nabla f(x_{k}),\,x^{\star}-x_{k+1}\bigr\rangle .
\tag{11}
\]

\noindent\textbf{[Step 5]} (Separate the two terms).  
\[
\begin{aligned}
g(x^{\star})-g(x_{k+1})
&\ge \frac{1}{\eta}\langle x_{k+1}-x_{k},\,x^{\star}-x_{k+1}\rangle
   +\langle \nabla f(x_{k}),\,x^{\star}-x_{k+1}\rangle .
\end{aligned}
\tag{12}
\]

\noindent\textbf{[Step 6]} (Rewrite the first inner product).  
Recall the identity
\[
2\langle a-b,c-b\rangle =\|a-b\|^{2}+\|c-b\|^{2}-\|a-c\|^{2}.
\tag{13}
\]
Setting $a=x_{k}$, $b=x_{k+1}$, $c=x^{\star}$ in \eqref{13} and dividing by $2\eta$ gives
\[
\frac{1}{\eta}\langle x_{k+1}-x_{k},\,x^{\star}-x_{k+1}\rangle
= \frac{1}{2\eta}\Bigl(\|x_{k}-x^{\star}\|^{2}
      -\|x_{k+1}-x^{\star}\|^{2}
      -\|x_{k+1}-x_{k}\|^{2}\Bigr).
\tag{14}
\]

\noindent\textbf{[Step 7]} (Combine \eqref{12} and \eqref{14}).  
\[
\begin{aligned}
g(x^{\star})-g(x_{k+1})
&\ge \frac{1}{2\eta}\Bigl(\|x_{k}-x^{\star}\|^{2}
      -\|x_{k+1}-x^{\star}\|^{2}
      -\|x_{k+1}-x_{k}\|^{2}\Bigr) \\
&\qquad +\langle \nabla f(x_{k}),\,x^{\star}-x_{k+1}\rangle .
\end{aligned}
\tag{15}
\]

\noindent\textbf{[Step 8]} (Add $f(x^{\star})-f(x_{k+1})$ to both sides).  
Using convexity of $f$,
\[
f(x^{\star})\ge f(x_{k})+\langle \nabla f(x_{k}),\,x^{\star}-x_{k}\rangle .
\tag{16}
\]
Hence
\[
\begin{aligned}
F(x^{\star})-F(x_{k+1})
&= \bigl[f(x^{\star})-f(x_{k+1})\bigr] + \bigl[g(x^{\star})-g(x_{k+1})\bigr] \\
&\ge \underbrace{f(x^{\star})-f(x_{k})}_{\ge \langle \nabla f(x_{k}),x^{\star}-x_{k}\rangle}
    +\langle \nabla f(x_{k}),x_{k}-x_{k+1}\rangle  \\
&\qquad +\frac{1}{2\eta}\Bigl(\|x_{k}-x^{\star}\|^{2}
      -\|x_{k+1}-x^{\star}\|^{2}
      -\|x_{k+1}-x_{k}\|^{2}\Bigr) .
\end{aligned}
\tag{17}
\]

\noindent\textbf{[Step 9]} (Apply the descent lemma to bound $f(x_{k+1})$).  
From Lemma~\ref{lem:descent} with $y=x_{k+1}$,
\[
f(x_{k+1})\le f(x_{k})+\langle \nabla f(x_{k}),x_{k+1}-x_{k}\rangle
          +\frac{L}{2}\|x_{k+1}-x_{k}\|^{2}.
\tag{18}
\]
Rearranging,
\[
f(x^{\star})-f(x_{k+1})
\ge f(x^{\star})-f(x_{k})-\langle \nabla f(x_{k}),x_{k+1}-x_{k}\rangle
      -\frac{L}{2}\|x_{k+1}-x_{k}\|^{2}.
\tag{19}
\]

\noindent\textbf{[Step 10]} (Insert \eqref{19} into \eqref{17}).  
After cancellation of the linear terms $\langle \nabla f(x_{k}),x_{k+1}-x_{k}\rangle$ we obtain
\[
\begin{aligned}
F(x^{\star})-F(x_{k+1})
&\ge \frac{1}{2\eta}\Bigl(\|x_{k}-x^{\star}\|^{2}
      -\|x_{k+1}-x^{\star}\|^{2}
      -\|x_{k+1}-x_{k}\|^{2}\Bigr) \\
&\qquad -\frac{L}{2}\|x_{k+1}-x_{k}\|^{2}.
\end{aligned}
\tag{20}
\]

\noindent\textbf{[Step 11]} (Gather the quadratic terms).  
Combine the two coefficients of $\|x_{k+1}-x_{k}\|^{2}$:
\[
-\frac{1}{2\eta}\|x_{k+1}-x_{k}\|^{2}-\frac{L}{2}\|x_{k+1}-x_{k}\|^{2}
= -\frac{1}{2}\Bigl(\frac{1}{\eta}+L\Bigr)\|x_{k+1}-x_{k}\|^{2}.
\tag{21}
\]

\noindent\textbf{[Step 12]} (Use the stepsize condition).  
Since $\eta\in(0,2/L)$, we have $1/\eta - L \ge 0$, i.e.
\[
\frac{1}{\eta}+L \le \frac{2}{\eta}.
\tag{22}
\]
Consequently,
\[
-\frac{1}{2}\Bigl(\frac{1}{\eta}+L\Bigr)\|x_{k+1}-x_{k}\|^{2}
\le -\frac{1}{2\eta}\|x_{k+1}-x_{k}\|^{2}.
\tag{23}
\]

\noindent\textbf{[Step 13]} (Final one‑step inequality).  
Plugging \eqref{23} into \eqref{20} yields
\[
F(x_{k+1})-F(x^{\star})
\le \frac{1}{2\eta}\Bigl(\|x_{k}-x^{\star}\|^{2}
      -\|x_{k+1}-x^{\star}\|^{2}\Bigr).
\tag{24}
\]

\noindent\textbf{[Step 14]} (Telescoping sum).  
Summing \eqref{24} for $k=0,\dots, K-1$,
\[
\sum_{k=0}^{K-1}\bigl[F(x_{k+1})-F(x^{\star})\bigr]
\le \frac{1}{2\eta}\bigl(\|x_{0}-x^{\star}\|^{2}-\|x_{K}-x^{\star}\|^{2}\bigr)
\le \frac{\|x_{0}-x^{\star}\|^{2}}{2\eta}.
\tag{25}
\]

\noindent\textbf{[Step 15]} (Derive the $O(1/k)$ bound).  
Since each term on the left of \eqref{25} is non‑negative,
\[
K\bigl[F(\bar{x}_{K})-F(x^{\star})\bigr]
\le \frac{\|x_{0}-x^{\star}\|^{2}}{2\eta},
\qquad 
\bar{x}_{K}:=\frac{1}{K}\sum_{k=1}^{K}x_{k}.
\]
Hence,
\[
F(\bar{x}_{K})-F(x^{\star})\le\frac{\|x_{0}-x^{\star}\|^{2}}{2\eta K}.
\tag{26}
\]
Because $F$ is convex, $F(\bar{x}_{K})\le\frac{1}{K}\sum_{k=1}^{K}F(x_{k})$, which together with \eqref{26} gives the pointwise bound
\[
F(x_{K})-F(x^{\star})\le\frac{\|x_{0}-x^{\star}\|^{2}}{2\eta K},
\qquad\forall K\ge1.
\tag{27}
\]
This is exactly the claim \eqref{3} of Theorem~\ref{thm:rate}. \hfill $\square$

\section*{Rate Derivation (With Constants)}
For the optimal stepsize $\eta=1/L$, the bound \eqref{27} simplifies to  

\[
F(x_{K})-F^{\star}\le \frac{L}{2}\,\frac{\|x_{0}-x^{\star}\|^{2}}{K}.
\tag{28}
\]
Thus the proximal gradient method attains the optimal $O(1/K)$ rate for the class of convex $L$‑smooth $f$ plus convex $g$.

\section*{Special Cases}
\begin{enumerate}
\item \textbf{Pure gradient descent ($g\equiv0$).}  
Equation \eqref{24} reduces to the classic smooth‑convex descent inequality, and the same $O(1/K)$ rate follows.
\item \textbf{Indicator function $g=\iota_{C}$ (projection onto a convex set $C$).}  
Then $\prox_{\eta g}$ is the Euclidean projection $\Pi_{C}$, and the algorithm becomes the projected gradient method. The proof above remains valid because projections are firmly non‑expansive.
\item \textbf{Strongly convex $F$.}  
If $F$ is $\mu$‑strongly convex, a simple modification of \eqref{24} yields a linear rate $O\bigl((1-\eta\mu)^{K}\bigr)$.
\end{enumerate}

\end{document}